% Aufgabentypen für Statistik nach Kapitel sortiert
% Basierend auf der Struktur aus StatistikUebungen.pdf und den Lösungsdateien
% Format orientiert an der Probeklausur

\documentclass[12pt, a4paper, titlepage, ngerman]{scrartcl}

\usepackage[utf8]{inputenc}
\usepackage[ngerman]{babel}
\usepackage{amsmath}
\usepackage{amssymb}
\usepackage{geometry}
\usepackage{fancyhdr}
\usepackage{listings}
\usepackage{xcolor}
\usepackage{booktabs}
\usepackage{longtable}

\geometry{left=2.5cm, right=2.5cm, top=2.5cm, bottom=2.5cm}

\pagestyle{fancy}
\fancyhf{}
\fancyhead[L]{Aufgabentypen -- Statistik}
\fancyhead[R]{\thepage}
\renewcommand{\headrulewidth}{0.4pt}

\title{Aufgabentypen Statistik}
\subtitle{Übersicht der Hauptaufgabentypen nach Kapitel}
\author{}
\date{}

\begin{document}

\maketitle

\tableofcontents

\newpage

% ============================================================================
\section{Übungen zu Wahrscheinlichkeitsrechnung und Zufallsvariablen}
% ============================================================================

\subsection{Reviewfragen}

\textbf{Aufgabentyp:} Grundlegende konzeptuelle Fragen zur Wahrscheinlichkeitsrechnung

\begin{description}
  \item[Gegebenes:] Aussagen oder Definitionen zu wahrscheinlichkeitstheoretischen Konzepten
  \item[Gesucht:] Beurteilung der Aussagen (wahr/falsch) oder Ergänzung von Definitionen
  \item[Methode:] Theoretisches Wissen anwenden, Unterscheidungen treffen (z.B. Elementarereignisse vs. Ereignisse)
\end{description}

\medskip

\textbf{Beispiel-Schema:}
\begin{itemize}
  \item (a) Elementarereignisse schließen sich gegenseitig aus, während Ereignisse sich nicht gegenseitig ausschließen.
  \item (b) Alle sind disjunkt.
  \item (c) Alle sind gleichwahrscheinlich.
\end{itemize}

\subsection{Laplace-Wahrscheinlichkeit}

\textbf{Aufgabentyp:} Berechnung von Wahrscheinlichkeiten bei endlichen Grundmengen mit gleichwahrscheinlichen Elementarereignissen

\begin{description}
  \item[Gegebenes:] Anzahl möglicher und günstiger Fälle (häufig bei Würfeln, Karten, etc.)
  \item[Gesucht:] $P(A) = \frac{\text{Anzahl günstiger Fälle}}{\text{Anzahl möglicher Fälle}}$
  \item[Methode:] Kombinatorik verwenden, verschiedene Ereignisse betrachten
\end{description}

\medskip

\textbf{Beispiel-Schema:}
\begin{itemize}
  \item (a) Die Anzahl der möglichen Würfe ist \ldots. Anzahl der günstigen Fälle für das Ereignis \ldots ist \ldots
  \item (b) Berechne $P(A)$ für verschiedene Szenarien
\end{itemize}

\subsection{Urnenmodell}

\textbf{Aufgabentyp:} Kombinatorische Probleme mit Variation und Kombination (mit/ohne Zurücklegen)

\begin{description}
  \item[Gegebenes:] Urne mit verschiedenen Objekten, Anzahl Züge, mit/ohne Zurücklegen
  \item[Gesucht:] Wahrscheinlichkeiten verschiedener Ereignisse
  \item[Methode:] 
    \begin{itemize}
      \item Variation ohne Wiederholung: $V(n,k) = \frac{n!}{(n-k)!}$
      \item Variation mit Wiederholung: $V_w(n,k) = n^k$
      \item Kombination ohne Wiederholung: $K(n,k) = \binom{n}{k}$
      \item Kombination mit Wiederholung: $K_w(n,k) = \binom{n+k-1}{k}$
    \end{itemize}
\end{description}

\subsection{Rechenregeln für Wahrscheinlichkeiten}

\textbf{Aufgabentyp:} Anwendung von Axiomen und Sätzen der Wahrscheinlichkeit

\begin{description}
  \item[Gegebenes:] Wahrscheinlichkeiten einzelner oder kombinierter Ereignisse
  \item[Gesucht:] Wahrscheinlichkeit eines zusammengesetzten Ereignisses
  \item[Methode:] Verwendung von:
    \begin{itemize}
      \item Komplementregel: $P(\overline{A}) = 1 - P(A)$
      \item Additionssatz: $P(A \cup B) = P(A) + P(B) - P(A \cap B)$
      \item Multiplikationssatz für unabhängige Ereignisse: $P(A \cap B) = P(A) \cdot P(B)$
    \end{itemize}
\end{description}

\subsection{Bedingte Wahrscheinlichkeiten}

\textbf{Aufgabentyp:} Berechnung bedingter Wahrscheinlichkeiten, Satz von Bayes

\begin{description}
  \item[Gegebenes:] Wahrscheinlichkeiten verschiedener Ereignisse, bedingte Wahrscheinlichkeiten
  \item[Gesucht:] $P(A|B)$ oder Umkehrung: $P(B|A)$
  \item[Methode:] 
    \begin{itemize}
      \item Bedingte Wahrscheinlichkeit: $P(A|B) = \frac{P(A \cap B)}{P(B)}$
      \item Satz von Bayes: $P(A|B) = \frac{P(B|A) \cdot P(A)}{P(B|A) \cdot P(A) + P(B|\overline{A}) \cdot P(\overline{A})}$
      \item Satz der totalen Wahrscheinlichkeit
    \end{itemize}
\end{description}

\subsection{Zufallsvariablen}

\textbf{Aufgabentyp:} Definition und Berechnung von Wahrscheinlichkeitsfunktionen

\begin{description}
  \item[Gegebenes:] Zufallsexperiment, mögliche Ausgänge und deren Wahrscheinlichkeiten
  \item[Gesucht:] Wahrscheinlichkeitsfunktion $f(x)$, Verteilungsfunktion $F(x)$, oder spezifische Wahrscheinlichkeiten
  \item[Methode:] 
    \begin{itemize}
      \item $f(x_i) = P(X = x_i)$ für diskrete Zufallsvariablen
      \item $F(x) = P(X \leq x) = \sum_{x_i \leq x} f(x_i)$
      \item $\int_{-\infty}^{\infty} f(x) dx = 1$ für stetige Zufallsvariablen
    \end{itemize}
\end{description}

% ============================================================================
\section{Übungen zu Verteilungen}
% ============================================================================

\subsection{Reviewfragen}

\textbf{Aufgabentyp:} Konzeptuelle Fragen zu verschiedenen Verteilungen

\subsection{Spezielle diskrete Verteilungen (Bowling, Würfel, Schachfiguren, Fußball)}

\textbf{Aufgabentyp:} Anwendung diskreter Verteilungen auf konkrete praktische Probleme

\begin{description}
  \item[Gegebenes:] Konkretes Szenario mit diskretem Merkmalsraum
  \item[Gesucht:] Wahrscheinlichkeiten, Erwartungswert, Varianz
  \item[Methode:] Identifikation der Verteilung (z.B. Binomial, Poisson), Anwendung der Verteilungsparameter
\end{description}

\subsection{Gleichverteilung}

\textbf{Aufgabentyp:} Stetige Gleichverteilung auf einem Intervall

\begin{description}
  \item[Gegebenes:] Intervall $[a, b]$
  \item[Gesucht:] Wahrscheinlichkeiten, Erwartungswert $E(X) = \frac{a+b}{2}$, Varianz $\text{Var}(X) = \frac{(b-a)^2}{12}$
  \item[Dichtefunktion:] $f(x) = \begin{cases} \frac{1}{b-a} & \text{für } a \leq x \leq b \\ 0 & \text{sonst} \end{cases}$
\end{description}

\subsection{Exponentialverteilung (Warteschlange)}

\textbf{Aufgabentyp:} Exponentialverteilung für Wartezeiten

\begin{description}
  \item[Gegebenes:] Parameter $\lambda > 0$, Wahrscheinlichkeit $P(X \leq t)$
  \item[Gesucht:] Parameter $\lambda$, oder Wahrscheinlichkeiten
  \item[Methode:] 
    \begin{itemize}
      \item Verteilungsfunktion: $F(t) = 1 - e^{-\lambda t}$
      \item Dichtefunktion: $f(t) = \lambda e^{-\lambda t}$
      \item Auflösen: $P(X \leq t) = p \Rightarrow 1 - e^{-\lambda t} = p \Rightarrow \lambda = -\frac{\ln(1-p)}{t}$
    \end{itemize}
\end{description}

\subsection{Normalverteilung}

\textbf{Aufgabentyp:} Anwendung der Normalverteilung (standardisiert und nicht-standardisiert)

\begin{description}
  \item[Gegebenes:] Parameter $\mu$ und $\sigma$, oder standardnormalverteilte Werte
  \item[Gesucht:] Wahrscheinlichkeiten, Werte für gegebene Wahrscheinlichkeiten
  \item[Methode:] 
    \begin{itemize}
      \item Standardisierung: $Z = \frac{X - \mu}{\sigma}$
      \item Lookup in Standardnormalverteilungs-Tabelle
      \item Eigenschaften der Normalverteilung nutzen
    \end{itemize}
\end{description}

\subsection{Drei-Sigma-Bereich}

\textbf{Aufgabentyp:} Anwendung der Drei-Sigma-Regel

\begin{description}
  \item[Gegebenes:] Normalverteilte Zufallsvariable
  \item[Gesucht:] Wahrscheinlichkeit im Bereich $[\mu - 3\sigma, \mu + 3\sigma]$
  \item[Methode:] $P(\mu - 3\sigma \leq X \leq \mu + 3\sigma) \approx 0.9974$
\end{description}

% ============================================================================
\section{Übungen zu Verteilungsparameter, Mehrdimensionale Zufallsvariablen}
% ============================================================================

\subsection{Lageparameter}

\textbf{Aufgabentyp:} Berechnung von Erwartungswert, Varianz und Median

\begin{description}
  \item[Gegebenes:] Wahrscheinlichkeitsfunktion oder Dichtefunktion einer Zufallsvariablen
  \item[Gesucht:] $E(X)$, $\text{Var}(X) = E((X - E(X))^2)$, Median
  \item[Methode:] 
    \begin{itemize}
      \item Diskret: $E(X) = \sum x_i \cdot f(x_i)$, $\text{Var}(X) = E(X^2) - (E(X))^2$
      \item Stetig: $E(X) = \int x \cdot f(x) dx$
      \item Median: kleinstes $x$ mit $F(x) \geq 0.5$
    \end{itemize}
\end{description}

\subsection{Gemeinsame Wahrscheinlichkeitsfunktion}

\textbf{Aufgabentyp:} Berechnung von Randwahrscheinlichkeiten und Abhängigkeit

\begin{description}
  \item[Gegebenes:] Gemeinsame Wahrscheinlichkeitsfunktion $f(x,y) = P(X = x \text{ und } Y = y)$
  \item[Gesucht:] Randwahrscheinlichkeiten $f_X(x)$, $f_Y(y)$, Abhängigkeit prüfen
  \item[Methode:] 
    \begin{itemize}
      \item Randwahrscheinlichkeit: $f_X(x) = \sum_y f(x,y)$
      \item Unabhängigkeit: $f(x,y) = f_X(x) \cdot f_Y(y)$ für alle $(x,y)$
    \end{itemize}
\end{description}

\subsection{Zweidimensionale diskrete/stetige Zufallsvariablen}

\textbf{Aufgabentyp:} Berechnung von Wahrscheinlichkeiten und Erwartungswerten für mehrdimensionale Variablen

\begin{description}
  \item[Gegebenes:] Gemeinsame Verteilung zweier Zufallsvariablen
  \item[Gesucht:] Wahrscheinlichkeiten, Erwartungswerte, Kovarianzen
  \item[Methode:] Integration oder Summation über den relevanten Bereich
\end{description}

\subsection{Erwartungswert}

\textbf{Aufgabentyp:} Berechnung von Erwartungswerten für Funktionen von Zufallsvariablen

\begin{description}
  \item[Gegebenes:] Zufallsvariable $X$ mit bekannter Verteilung, Funktion $g(X)$
  \item[Gesucht:] $E(g(X))$
  \item[Methode:] $E(g(X)) = \sum g(x_i) f(x_i)$ oder $E(g(X)) = \int g(x) f(x) dx$
\end{description}

% ============================================================================
\section{Übungen zu Grundlagen der induktiven Statistik}
% ============================================================================

\subsection{Likelihoodfunktion}

\textbf{Aufgabentyp:} Berechnung und Analyse der Likelihoodfunktion

\begin{description}
  \item[Gegebenes:] Stichprobe, Verteilungsannahme
  \item[Gesucht:] Likelihoodfunktion $L(\theta) = \prod_{i=1}^n f(x_i | \theta)$
  \item[Methode:] Produkt der Wahrscheinlichkeits- oder Dichtefunktionen bilden
\end{description}

\subsection{Erwartungswert und Varianz von Stichprobenfunktionen}

\textbf{Aufgabentyp:} Berechnung von Eigenschaften von Schätzfunktionen

\begin{description}
  \item[Gegebenes:] Grundgesamtheit mit Erwartungswert $\mu$ und Varianz $\sigma^2$, Stichprobengröße $n$
  \item[Gesucht:] $E(\overline{X})$, $\text{Var}(\overline{X})$
  \item[Methode:] 
    \begin{itemize}
      \item $E(\overline{X}) = \mu$
      \item $\text{Var}(\overline{X}) = \frac{\sigma^2}{n}$
      \item Standardfehler: $\sigma_{\overline{X}} = \frac{\sigma}{\sqrt{n}}$
    \end{itemize}
\end{description}

\subsection{Fraktile}

\textbf{Aufgabentyp:} Bestimmung von Quantilen einer Verteilung

\begin{description}
  \item[Gegebenes:] Wahrscheinlichkeit $p$, Verteilung
  \item[Gesucht:] $x_p$ mit $P(X \leq x_p) = p$
  \item[Methode:] Umkehrung der Verteilungsfunktion $F^{-1}(p)$
\end{description}

% ============================================================================
\section{Aufgaben zur Punktschätzung}
% ============================================================================

\subsection{Erwartungstreue und Wirksamkeit}

\textbf{Aufgabentyp:} Prüfung und Vergleich von Schätzfunktionen

\begin{description}
  \item[Gegebenes:] Schätzfunktion $\hat{\theta}$
  \item[Gesucht:] 
    \begin{itemize}
      \item Erwartungstreue: $E(\hat{\theta}) = \theta$?
      \item Wirksamkeit (Effizienz): welche Schätzfunktion hat kleinere Varianz?
    \end{itemize}
  \item[Methode:] Berechnung von $E(\hat{\theta})$ und $\text{Var}(\hat{\theta})$
\end{description}

\subsection{Maximum-Likelihood-Schätzung}

\textbf{Aufgabentyp:} Bestimmung des Maximum-Likelihood-Schätzers

\begin{description}
  \item[Gegebenes:] Likelihoodfunktion $L(\theta)$
  \item[Gesucht:] Parameter $\hat{\theta}_{ML}$ der Likelihoodfunktion maximiert
  \item[Methode:] 
    \begin{itemize}
      \item Logarithmische Likelihoodfunktion: $\ell(\theta) = \ln L(\theta)$
      \item Erste Ableitung Null setzen: $\frac{d\ell(\theta)}{d\theta} = 0$
      \item Auflösen nach $\theta$
    \end{itemize}
\end{description}

\subsection{Bayes-Schätzung}

\textbf{Aufgabentyp:} Schätzung unter Verwendung von Priori-Information

\begin{description}
  \item[Gegebenes:] Priori-Verteilung, Likelihoodfunktion, Daten
  \item[Gesucht:] Posteriori-Schätzung des Parameters
  \item[Methode:] Satz von Bayes zur Kombination von Priori und Likelihood
\end{description}

\subsection{Konsistenz}

\textbf{Aufgabentyp:} Überprüfung, ob Schätzfunktion konsistent ist

\begin{description}
  \item[Gegebenes:] Schätzfunktion $\hat{\theta}_n$
  \item[Gesucht:] Konvergiert $\hat{\theta}_n \to \theta$ für $n \to \infty$?
  \item[Methode:] Tschebyscheff-Ungleichung oder andere Konvergenzkriterien
\end{description}

% ============================================================================
\section{Übungen zur Intervallschätzung}
% ============================================================================

\subsection{Schätzintervall (Konfidenzintervall)}

\textbf{Aufgabentyp:} Berechnung von Konfidenzintervallen für Parameter

\begin{description}
  \item[Gegebenes:] Stichprobe, Konfidenzniveau $1 - \alpha$
  \item[Gesucht:] Konfidenzintervall $[\text{untere Grenze}, \text{obere Grenze}]$
  \item[Schema:] 
    \begin{enumerate}
      \item Konfidenzniveau festlegen: $1 - \alpha$
      \item Kritischen Wert bestimmen (z.B. $z_{\alpha/2}$ aus Standard-Normalverteilung oder $t_{\alpha/2}$ aus t-Verteilung)
      \item Konfidenzintervall berechnen: $\text{Schätzer} \pm \text{kritischer Wert} \cdot \text{Standardfehler}$
    \end{enumerate}
\end{description}

\subsubsection{Symmetrisches Konfidenzintervall bei dichotomer Verteilung}

\textbf{Aufgabentyp:} Konfidenzintervall für einen Anteilswert

\begin{description}
  \item[Gegebenes:] Anzahl der Erfolge $m$ von $n$ Versuchen, Konfidenzniveau $1 - \alpha$
  \item[Gesucht:] Konfidenzintervall $[d_u, d_o]$ für Anteil $p$
  \item[Voraussetzung:] $5 \leq m \leq n - 5$ (Approximation mit Normalverteilung)
  \item[Schema:] 
    \begin{enumerate}
      \item Punktschätzer: $\hat{p} = \frac{m}{n}$
      \item Standardfehler: $\hat{\sigma} = \sqrt{\hat{p}(1 - \hat{p})}$
      \item Kritischer Wert: $c = z_{1 - \alpha/2}$
      \item Konfidenzintervall: $\left[\hat{p} - \frac{\hat{\sigma} \cdot c}{\sqrt{n}}, \hat{p} + \frac{\hat{\sigma} \cdot c}{\sqrt{n}}\right]$
    \end{enumerate}
\end{description}

\subsubsection{Konfidenzintervall für Mittelwert (Normalverteilung)}

\textbf{Aufgabentyp:} Konfidenzintervall für $\mu$ bei bekannter oder unbekannter Varianz

\begin{description}
  \item[Gegebenes:] Stichprobe $x_1, \ldots, x_n$, Konfidenzniveau
  \item[Gesucht:] Konfidenzintervall $[d_u, d_o]$ für $\mu$
  \item[Methode:] 
    \begin{itemize}
      \item $\sigma$ bekannt: $\overline{x} \pm z_{\alpha/2} \cdot \frac{\sigma}{\sqrt{n}}$
      \item $\sigma$ unbekannt: $\overline{x} \pm t_{\alpha/2, n-1} \cdot \frac{s}{\sqrt{n}}$ (t-Test)
    \end{itemize}
\end{description}

% ============================================================================
\section{Übungen zu Signifikanztests}
% ============================================================================

\subsection{Einstichproben-t-Test}

\textbf{Aufgabentyp:} Hypothesentest für Mittelwert bei unbekannter Varianz

\begin{description}
  \item[Voraussetzungen:] $n < 30$, $X \sim N(\mu, \sigma)$, $\sigma$ unbekannt
  \item[Hypothesen:] 
    \begin{itemize}
      \item Rechtsseitig: $H_0: \mu \leq \mu_0$ vs. $H_1: \mu > \mu_0$
      \item Linksseitig: $H_0: \mu \geq \mu_0$ vs. $H_1: \mu < \mu_0$
      \item Zweiseitig: $H_0: \mu = \mu_0$ vs. $H_1: \mu \neq \mu_0$
    \end{itemize}
  \item[Testgröße:] $t = \frac{\overline{x} - \mu_0}{\frac{s}{\sqrt{n}}} \sim t(n-1)$
  \item[Schema:] 
    \begin{enumerate}
      \item Hypothesen formulieren
      \item Signifikanzniveau $\alpha$ festlegen
      \item Stichprobenmittelwert und Stichprobenvarianz berechnen
      \item Testgröße berechnen
      \item Kritischen Wert bestimmen (t-Tabelle)
      \item Entscheidung treffen
    \end{enumerate}
\end{description}

\subsection{Chi-Quadrat-Anpassungstest}

\textbf{Aufgabentyp:} Hypothesentest auf Goodness of Fit

\begin{description}
  \item[Gegebenes:] Beobachtete Häufigkeitsverteilung, hypothetische Verteilung
  \item[Gesucht:] Passt die Beobachtung zu der hypothetischen Verteilung?
  \item[Hypothesen:] $H_0: F = F_0$ (hypothetische Verteilung) vs. $H_1: F \neq F_0$
  \item[Testgröße:] $\chi^2 = \sum_{i=1}^k \frac{(h_i - np_i)^2}{np_i} \sim \chi^2(k-1)$
    \begin{itemize}
      \item $h_i$ = beobachtete Häufigkeit
      \item $np_i$ = erwartete Häufigkeit
      \item Voraussetzung: $\forall i: 5 < np_i$
    \end{itemize}
  \item[Schema:] 
    \begin{enumerate}
      \item Hypothesen formulieren
      \item Signifikanzniveau $\alpha$ festlegen
      \item Erwartete Häufigkeiten berechnen
      \item Testgröße berechnen
      \item Kritischen Wert: $\chi^2_{\alpha}(k-1)$ bestimmen
      \item Entscheidung: Wenn $\chi^2 > \chi^2_{\alpha}$, dann $H_0$ ablehnen
    \end{enumerate}
\end{description}

\subsection{Zweistichproben-t-Test}

\textbf{Aufgabentyp:} Vergleich von Mittelwerten zweier Stichproben

\begin{description}
  \item[Voraussetzungen:] Beide Stichproben normalverteilt, Varianzen gleich
  \item[Hypothesen:] 
    \begin{itemize}
      \item $H_0: \mu_1 = \mu_2$ vs. $H_1: \mu_1 \neq \mu_2$ (zweiseitig)
      \item $H_0: \mu_1 \leq \mu_2$ vs. $H_1: \mu_1 > \mu_2$ (einseitig)
    \end{itemize}
  \item[Testgröße:] $t = \frac{\overline{x}_1 - \overline{x}_2}{s_p \sqrt{\frac{1}{n_1} + \frac{1}{n_2}}} \sim t(n_1 + n_2 - 2)$
    \begin{itemize}
      \item Gepoolte Varianz: $s_p^2 = \frac{(n_1-1)s_1^2 + (n_2-1)s_2^2}{n_1 + n_2 - 2}$
    \end{itemize}
\end{description}

\subsection{Signifikanztest (allgemein)}

\textbf{Aufgabentyp:} Hypothesentest für verschiedene Parameter

\begin{description}
  \item[Allgemeines Schema:] 
    \begin{enumerate}
      \item Hypothesen aufstellen: $H_0$ (Nullhypothese) und $H_1$ (Alternativhypothese)
      \item Signifikanzniveau $\alpha$ wählen
      \item Teststatistik/Testgröße berechnen
      \item Ablehnungsbereich bzw. kritischen Wert bestimmen
      \item Entscheidung: 
        \begin{itemize}
          \item Wenn Testgröße im Ablehnungsbereich liegt: $H_0$ ablehnen
          \item Sonst: $H_0$ nicht ablehnen
        \end{itemize}
      \item Antwort im Sachkontext formulieren
    \end{enumerate}
\end{description}

\newpage

\appendix

\section{Wichtige Formeln}

\subsection{Wahrscheinlichkeitsrechnung}

\begin{itemize}
  \item Laplace-Wahrscheinlichkeit: $P(A) = \frac{|A|}{|\Omega|}$
  \item Bedingte Wahrscheinlichkeit: $P(A|B) = \frac{P(A \cap B)}{P(B)}$
  \item Satz von Bayes: $P(A|B) = \frac{P(B|A) \cdot P(A)}{P(B)}$
\end{itemize}

\subsection{Verteilungsparameter}

\begin{itemize}
  \item Erwartungswert (diskret): $E(X) = \sum x_i \cdot f(x_i)$
  \item Varianz: $\text{Var}(X) = E((X - E(X))^2) = E(X^2) - (E(X))^2$
  \item Standardabweichung: $\sigma = \sqrt{\text{Var}(X)}$
\end{itemize}

\subsection{Standardnormalverteilung}

\begin{itemize}
  \item Dichtefunktion: $\phi(z) = \frac{1}{\sqrt{2\pi}} e^{-z^2/2}$
  \item Verteilungsfunktion: $\Phi(z) = P(Z \leq z)$
\end{itemize}

\subsection{Stichprobenverteilungen}

\begin{itemize}
  \item Erwartungswert des Stichprobenmittels: $E(\overline{X}) = \mu$
  \item Varianz des Stichprobenmittels: $\text{Var}(\overline{X}) = \frac{\sigma^2}{n}$
  \item Standardfehler: $\sigma_{\overline{X}} = \frac{\sigma}{\sqrt{n}}$
\end{itemize}

\end{document}
